\subsection{Owner matchings and four-cycle obstructions}

Return to a physical bipartite partial template, and identify witnesses and
owners with $V$, $|V|=2b$. Thus $P$ joins equal physical colours, while $Q$
and the owner graph $B$ join opposite colours; $\Delta(B)\le2$.
The components of $B$ are paths and even cycles. Let $o$ be the number of
odd-order path components, including isolated vertices. Since $|V|$ is even,
$o$ is even.

Choose a matching $R$ on $V$ as follows. On an even-order path use its
perfect matching. On an odd-order path use either alternating matching that
leaves an endpoint uncovered. On a four-cycle pair the opposite vertices.
On every other even cycle use either perfect matching. These choices lie
inside their respective owner components. Let $W$ be the covered vertices,
so $|W|=2b-o$, and put $J_R=(P\cup Q\cup R)[W]$, with repeated edges
identified.

\begin{theorem}\label{thm:owner-matching}
If $J_R$ has $z$ components isomorphic to $K_4$, then
\[
 \beta_I(H)\ge
 b+\frac{2b+o}{2}+z+
       \left\lceil\frac{2b-o-4z}{3}\right\rceil
 =\left\lceil\frac{8b}{3}+\frac{o-2z}{6}\right\rceil.
\]
In particular, if some choice has $z=0$, then
$\beta_I(H)\ge\lceil8b/3+o/6\rceil$.
Such a choice always exists if $P\cup Q$ has no four-cycle component,
regardless of the lengths of the owner paths and even cycles.
\end{theorem}
\begin{proof}
The graph $J_R$ has maximum degree three. Every $K_4$ is an entire component,
since its vertices already have degree three. All its incident matching types
are distinct, so deleting its $R$ matching leaves a four-cycle in $P\cup Q$.
This four-cycle is a whole component even before restriction to $W$, because
both $P$ and $Q$ degrees are exhausted. Thus no such $K_4$ is possible when
$P\cup Q$ has no four-cycle component.

Every component other than $K_4$ is three-colourable by Brooks's theorem
\cite{brooks}; odd cycles also admit three colours, and lower-degree
components cause no exception. Take one witness index in each $K_4$ and a
largest colour class in their complement. The resulting independent set $T$
of $J_R$ has size at least
$z+\lceil(2b-o-4z)/3\rceil$.

Select all owners indexed by $V\setminus W$, and one endpoint outside $T$
from every edge of $R$, obtaining $S_C$. Then
$|S_C|=(2b+o)/2$ and $S_C\cap T=\varnothing$.
On every owner path and every non-square cycle, $B-R$ is a matching;
no $R$ edge has both endpoints in $S_C$. On each square exactly two owners
are chosen, so their induced graph also has maximum degree at most one.
There are no owner edges between components. Hence $B[S_C]$ is dissociation.
Also $T$ is independent in $P\cup Q$. All helpers, these owners and these
witnesses therefore form a dissociation set, by Theorem~\ref{thm:all-helper}.
Adding their cardinalities proves the bound. The rounding identity follows
because $o$ is even. Every matching edge, uncovered endpoint and helper is
used once.
\end{proof}

The three-colouring is an existence argument with the usual constructive
Brooks theorem available. The verification implementation uses a colouring
search and makes no polynomial runtime claim for that search.
The theorem covers, for example, owner eight-cycles and ten-cycles, including
mixtures with paths and squares, whenever the auxiliary four-cycle obstruction
is absent. It imposes a condition on the actual template, rather than on an
owner presentation alone.

\begin{corollary}\label{cor:owner-matching-gap}
If the actual fixed witness extraction of a minimal forming set $S$ in a
bipartite graph of maximum degree three admits a $K_4$-free choice above, then
$51(|S|-\diss(G))\le5|V(G)|$.
\end{corollary}
\begin{proof}
The template has $b'=b-b_3$ helpers. The retained set of size $L$ is disjoint
from it and has no edge to it. The theorem gives
$Z=3\diss(G)-3L-8b'\ge0$. Substitution in the same fixed-extraction identity
as in Corollary~\ref{cor:path-six-core} gives
\[
 5n-51d=Z+5r+12(\alpha+\pi)+24\tau+7(e-b_3)+2b_3+5q\ge0.
\]
\end{proof}
This also bounds $\SF-\diss$ if a maximum-cardinality minimal forming set
has such an extraction. The coefficient is the existing conditional $5/51$;
the new point is the additional class of templates.

\paragraph{An exact Boolean test for the specified choices.}
Assign one Boolean variable to each non-square cycle and each odd-order path
of order at least three, choosing its alternating matching. The even-order
paths and square diagonals are fixed. For each auxiliary four-cycle, consider
its two missing diagonals. It becomes a $K_4$ exactly when both diagonals
belong to the chosen $R$. Each diagonal is absent from all choices, fixed,
or present in exactly one phase of one variable. Thus the forbidden event
is impossible, unconditional, a single literal, or a conjunction of two
literals. Opposite literals of one variable give an impossible event.
Negating the remaining events produces clauses of length at most two,
including an empty clause for an unconditional obstruction. The conjunction
is satisfiable if and only if this specified matching family contains a
$K_4$-free choice. The implication-graph test decides this in polynomial time
and returns either phases or contradictory directed paths. The two diagonals
may belong to different owner components; both variables must then be kept.

Failure of this test does not rule out the cardinality bound, other owner
matchings on paths, a different all-helper construction, or an unrestricted
dissociation set. The following connected template makes the distinction
explicit even when every owner perfect matching is included in the test.

\paragraph{A twelve-index obstruction with exact restricted optima.}
Let $V=\{0,\ldots,11\}$, with physical colours $0$ on $0,\ldots,5$ and $1$
on $6,\ldots,11$. Put
\begin{align*}
 P={}&\{01,23,45,67,89,(10)(11)\},\\
 Q={}&\{08,19,2(11),3(10),46,57\},
\end{align*}
where a juxtaposition denotes an unordered edge and parenthesized numbers
separate multi-digit endpoints. Let $B$ be the cycle in the order
\[
 0,7,4,8,3,11,2,10,1,6,5,9,0.
\]
Its two perfect matchings are the alternate edges $R_0,R_1$ in this order.
The physical template has thirty vertices, is connected and bipartite, and
has minimum degree two and maximum degree three: every witness and owner
has degree three, and every helper has degree two.

For $R_0$, $J_R$ has the $K_4$ component $\{2,3,10,11\}$; its other eight
vertices have the spanning cycle
\[
 0,1,6,4,8,9,5,7,0.
\]
For $R_1$, the $K_4$ is $\{4,5,6,7\}$, and the other cycle is
\[
 0,1,10,11,2,3,8,9,0.
\]
In each eight-cycle, the edges $08$ and $19$ block its two alternating
four-sets. An independent four-set in an eight-cycle must be alternating,
so the independence number of either eight-vertex remainder is at most
three. The independent sets $\{0,2,4,9\}$ and $\{0,2,4,10\}$ respectively
show that
\[
 \alpha(J_{R_0})=\alpha(J_{R_1})=4.
\]
Thus neither choice is $K_4$-free, although both satisfy the proposed
independence threshold $\lceil|V|/3\rceil=4$.

The larger all-helper family has exact optimum seventeen. A witness is given
by owner indices $\{3,4,6,7,9,10,11\}$ and witness indices $\{1,2,5,8\}$,
together with all six helpers. For the upper bound, Theorem~\ref{thm:all-helper}
gives $\beta_I\le18$. Equality would require $U$ independent in $P\cup Q$
and $B[V\setminus U]$ dissociation. Enlarge $U$ to a maximal independent
set; this preserves the latter condition. The auxiliary graph consists of
three four-cycles, whose opposite-pair choices are
\[
 (\{0,9\},\{1,8\}),\quad
 (\{2,10\},\{3,11\}),\quad
 (\{4,7\},\{5,6\}).
\]
Every maximal independent set chooses one pair per component. In each case
below, its owner complement contains the displayed three-vertex path:
\[
\begin{array}{c|c@{\qquad}c|c}
000&8,3,11&001&8,3,11\\
010&6,1,10&011&7,4,8\\
100&6,5,9&101&7,0,9\\
110&6,5,9&111&7,0,9
\end{array}
\]
The middle entry of each triple is its centre. None of the eight complements
is dissociation, excluding eighteen and proving $\beta_I=17$.

For $r$ disjoint copies, all owner perfect-matching choices therefore have
independence number $4r$, and the construction retaining exactly one owner
per matching edge has optimum $16r$, whereas $\beta_I=17r$.
This family is disconnected for $r>1$. It proves an unbounded additive loss
for that construction, with all helpers still retained in the better set.
It refutes neither the full-cycle matching-selection conjecture nor the
general $8b/3$ bound. Both remain open, as does the sharp unrestricted
subcubic coefficient.
